% ECONOMICS_PROBLEM: {"id": "C62-53", "title": "Recursive equilibrium without an auxiliary continuation state", "jel_code": "C62", "source_id": "OP-0582"}
% !TeX program = xelatex
\documentclass[11pt,a4paper]{article}
\usepackage[margin=23mm]{geometry}
\usepackage{fontspec}
\setmainfont{Latin Modern Roman}
\usepackage{amsmath,amssymb,mathtools}
\usepackage{xcolor,enumitem,booktabs,longtable,array,xurl,hyperref}
\definecolor{muted}{HTML}{53636C}
\hypersetup{colorlinks=true,urlcolor=blue,linkcolor=blue}
\setlength{\parindent}{0pt}
\setlength{\parskip}{5pt}
\newcommand{\R}{\mathbb R}
\newcommand{\E}{\mathbb E}
\newcommand{\N}{\mathbb N}
\newcommand{\ind}{\mathbf 1}
\newcommand{\Var}{\operatorname{Var}}
\newcommand{\Cov}{\operatorname{Cov}}
\newcommand{\supp}{\operatorname{supp}}
\DeclareMathOperator*{\argmin}{arg\,min}
\DeclareMathOperator*{\argmax}{arg\,max}
\newcommand{\entrymeta}[3]{{\sffamily\small\color{muted}\textbf{\texttt{#1}}\quad|\quad #2\par #3\par}}
\newcommand{\Problemref}[1]{\texttt{#1}}
\newcommand{\JELref}[2]{\texttt{#2} (JEL \texttt{#1})}
\begin{document}
\section*{Recursive equilibrium without an auxiliary continuation state}
\textbf{Problem ID:} \texttt{C62-53}\quad \textbf{JEL:} \texttt{C62}\par
% BEGIN ECONOMICS STATEMENT
\label{problem:OP-0582}
\label{beyond:MA-01}
\entrymeta{MA-01}{Heterogeneous-agent macroeconomics}{Published historical question, open in 2009}
\subsection*{Definitions and assumptions}
Consider infinitely lived price-taking households with common discount factor $\beta\in(0,1)$ and increasing strictly concave utility $u(c)$. Aggregate productivity $z$ and idiosyncratic labor efficiency $e>0$ take finitely many values, with specified aggregate probabilities $q(z'\mid z)$ and conditional efficiency probabilities $q(e'\mid e,z,z')$. Their product specifies the joint transition; aggregate shocks do not depend on individual efficiency. Households hold capital $a\ge\underline a$; there are no securities contingent on idiosyncratic shocks. The cross-sectional distribution $\mu$ of $(a,e)$ has finite first moment. Write $K=\int a\,d\mu>0$, $L=\int e\,d\mu>0$. A differentiable concave constant-returns production function $F(z,K,L)$ and depreciation $\delta\in(0,1)$ give factor prices
\[
 R(z,\mu)=1-\delta+F_K(z,K,L),\qquad w(z,\mu)=F_L(z,K,L).
\]
A recursive competitive equilibrium on an invariant set of distributions consists of a value $V$, savings rule $g$, and distribution transition $\mathcal T$ satisfying
\[
 V(a,e,z,\mu)=\sup_{\substack{a'\ge\underline a\\Ra+we-a'>0}}\left\{u(Ra+we-a')+
 \beta\E[V(a',e',z',\mathcal T(z,\mu,z'))\mid e,z]\right\},
\]
with $g$ attaining the supremum. For Borel $A$ and efficiency set $E$, distribution consistency means
\[
 \mathcal T(z,\mu,z')(A\times E)=\int \ind_{A}(g(a,e,z,\mu))q(E\mid e,z,z')\,d\mu(a,e).
\]
Require $V$ to equal the actual expected lifetime-utility supremum over feasible adapted plans under the induced aggregate transition, with finite values; a Bellman fixed point alone is insufficient. Capital and goods markets clear: $K'=\int g\,d\mu$ and $\int c\,d\mu+K'=F(z,K,L)+(1-\delta)K$. Assume well-defined expectations and a continuum law of large numbers conditional on aggregate shocks.
\subsection*{Formal question}
Under precise existence conditions for the corresponding sequential competitive economy, does an equilibrium of this recursive form exist from each admissible initial $(z_0,\mu_0)$? If additional assumptions are required, characterize them or exhibit failure within the specified incomplete-market family. The recursive aggregate state is $(z,\mu)$; no separately carried continuation-value coordinate is permitted.
\subsection*{Scope and status}
The full distribution $\mu$ is retained. The question does not ask for exact compression to aggregate capital or finitely many moments. The verified published review describes the auxiliary coordinate as individual expected continuation values. This historical question is not presented as a verified unresolved theorem in 2026.
\subsection*{References for the source of the problem}
Jonathan Heathcote, Kjetil Storesletten, and Giovanni L.~Violante, \emph{Quantitative Macroeconomics with Heterogeneous Households}, NBER Working Paper~14768 (March 2009), \url{https://www.nber.org/system/files/working_papers/w14768/w14768.pdf}. Verified published version: \emph{Annual Review of Economics} 1 (2009), 319--354, Section~5.2, footnote~20 on p.~342; \url{https://doi.org/10.1146/annurev.economics.050708.142922}. Author-hosted text: \url{https://violante.economics.princeton.edu/sites/g/files/toruqf5621/files/documents/Quantitative%20Macroeconomics%20with%20Heterogeneous%20Households.pdf}.

\subsection*{Common conventions: Source mathematical conventions}

\textbf{Well-defined objects.}
All probability spaces and measurable variables are specified or inherited
from a local statistical model. Expectations used as finite targets require
the stated integrability. A decision rule or estimator is measurable with
respect to its available information; randomized rules may use an auxiliary
independent random variable. Norms without a subscript are Euclidean in
finite-dimensional spaces. An optimizer is requested only when attainment
is assured; otherwise the question concerns the optimal value and
$\varepsilon$-optimal admissible rules for every $\varepsilon>0$.

\textbf{Identification.}
A structural model $m\in\mathcal M$ implies an observable law
$P_m^{\rm obs}$ and target $\theta(m)$. The target is point identified on
$\mathcal M$ if
\[
 P_{m_1}^{\rm obs}=P_{m_2}^{\rm obs}
 \quad\Longrightarrow\quad\theta(m_1)=\theta(m_2)
 \qquad(m_1,m_2\in\mathcal M).
\]
When this fails, the sharp identified set is
\[
 \Theta(P;\mathcal M)=\{\theta(m):m\in\mathcal M,
                                  P_m^{\rm obs}=P\}.
\]
Specifying an intervention or estimand does not establish identification.
Positivity, exclusion, causal sufficiency, measurement restrictions, and
transport assumptions are substantive local inputs, never automatic
consequences of notation.

\textbf{Minimax risk and nontrivial inference.}
For a statistical experiment with data law $P^{(n)}$, target $\theta(P)$,
and nonnegative loss $\ell$, define
\[
 \mathcal R_n(\mathcal P,\ell)
   =\inf_{\widehat\theta_n}\sup_{P\in\mathcal P}
      \E_{P^{(n)}}\ell(\widehat\theta_n,\theta(P)).
\]
The infimum is over measurable estimators. A sharp rate is a sequence
$r_n$ for which $0<c\leq C<\infty$, independent of $n$, satisfies
$c r_n\leq\mathcal R_n\leq C r_n$ for all sufficiently large $n$.
Squared-error parametric risk is of order $n^{-1}$; the corresponding
root-mean-squared error is of order $n^{-1/2}$. These different conventions
are distinguished locally. A uniform asymptotic confidence region $C_n$
must satisfy
\[
 \liminf_{n\to\infty}\inf_{P\in\mathcal P}
     \Pr_{P^{(n)}}\{\theta(P)\in C_n\}\geq 1-\alpha,
 \qquad 0<\alpha<1,
\]
together with an informative diameter or expected-width requirement.
Coverage alone permits the vacuous whole parameter space. Testing questions
similarly impose both size control and power against a declared separated
alternative class.

\textbf{Binary causal baseline.}
When an entry explicitly invokes this baseline, one observes independent
copies of $O=(X,W,Y)$, with $W\in\{0,1\}$ and potential outcomes $Y(0),Y(1)$.
Assume consistency $Y=Y(W)$, no interference, conditional exchangeability
$(Y(0),Y(1))\perp W\mid X$, and overlap
$\epsilon\leq e(X)\leq1-\epsilon$ almost surely for a fixed
$0<\epsilon<1/2$, where
\[
 e(x)=\Pr(W=1\mid X=x),\qquad
 m_w(x)=\E[Y\mid W=w,X=x],\qquad
 \tau(x)=m_1(x)-m_0(x).
\]
These assumptions identify conditional mean effects almost everywhere
under the covariate law. Pointwise or uniform effects require appropriate
regularity and a specified support. Continuous treatment, longitudinal
data, adaptive experiments, and network interference require their own
assumptions in place of this baseline. Bounded outcomes, densities, and
smoothness are additional assumptions only where stated.

\textbf{Function classes.}
On $[0,1]^d$, $\mathcal H_d^s(L)$ is the isotropic Hölder ball of radius
$L>0$: put $k=\lceil s\rceil-1$ and $\gamma=s-k\in(0,1]$, bound derivatives
through order $k$, and require the order-$k$ derivatives to be
$\gamma$-Hölder with constant at most $L$. Thus integer orders use a
Lipschitz final derivative convention. The class
$\operatorname{BV}_d(L)$ consists of functions $f\in L^1((0,1)^d)$ whose
distributional derivative is a finite vector measure and for which
$\|f\|_{L^1}+|Df|((0,1)^d)\leq L$.
An $L^\infty$ bound or a particular pointwise version is imposed separately
when needed. Sparse and neural-network classes require a declared
dictionary or architecture and coefficient bounds.

An anisotropic H\"older class with declared block smoothness indices means
coordinatewise H\"older regularity in each block, uniformly in the other
blocks, with all corresponding derivative and H\"older bounds at most
the stated radius. Derivative orders and the integer-order convention in
each block are those just given. Mixed-derivative bounds are additional
restrictions only when explicitly stated.

\textbf{Decision and computational questions.}
Policy regret compares admissible feasible policies under the same law and
constraints. In computational entries, finite data and thresholds have
rational encodings; running time is measured in their total bit length.
Real-valued equilibrium search uses the explicitly declared algebraic or
FIXP encoding. An approximation ratio for maximization is written
$\operatorname{OPT}/\operatorname{ALG}$ and for minimization as
$\operatorname{ALG}/\operatorname{OPT}$, with zero optima and infeasible
instances handled locally. Historical approximation bounds are not
automatically current bounds.
% END ECONOMICS STATEMENT
\end{document}
